\section{An initial curvature profile proved from the exact last source step}
Let $\kappa$ be an actual arm's first missing weight, let
$\delta=E(\kappa)$, and write $A=z^{-\kappa}E$.
Then
\[
\begin{array}{ll}
A_1-A_0\ge0,&\kappa\ge4,\\
A_n-A_{n-1}\ge2^{n-2},&2\le n\le\kappa-3,\\
A_n-A_{n-1}\ge\delta\,2^{n-2},&2\le n\le\kappa-\delta-1.
\end{array}                                                   \tag{1}
\]
These are SOURCE CURVATURE lower bounds. They are not obtained by
differentiating a coefficientwise lower bound.

Use the exact three normalized source recurrences, with the old
first weight denoted $k$:
\[
 1+z[k-1]A\quad(t=k-1),\qquad
 1+[k]A\quad(t=k),\qquad
 [t]A+z^{t-k}\quad(t>k).
\]
At every index in (1), the windows are uncut.
In a stationary reset the new increment at $n\ge2$ is exactly
the old coefficient $A_{n-1}$, bounded below by $2^{n-2}$.
The new leading multiplicity is one, so this also proves its scaled bound.
In a preserving or tied growth step the increment at $n\ge2$ is
the old coefficient $A_n$, plus a current unit birth and minus
a preceding unit birth. The unscaled early dyadic theorem gives
$A_n\ge2\,2^{n-2}$, paying that possible unit loss.
For the scaled claim a preserving step uses
$A_n\ge2\delta\,2^{n-2}$; a tied step uses
$2\delta_{\rm old}\ge\delta_{\rm old}+1=\delta_{\rm new}$.
The displayed rank limits put the old coefficient within exactly
the already proved dyadic or scaled-dyadic range.
The case $n=1$ follows separately: the three possible new
increments are $\delta_{\rm old}-1$, $A_1^{\rm old}-1$, or
$A_1^{\rm old}$ plus a nonnegative birth. The earlier bounds give
their nonnegativity when the new first weight is at least four.
The initial all-ones case has $\kappa=2$, so contributes no
nonempty range. This completes the unbounded last-step induction.

In particular EVERY actual arm of first missing weight at least
five has $\Delta E(\kappa+2)\ge1$. This gives a second compulsory
source-payment term in addition to the first birth.
The unscaled range in (1) cannot simply be extended by one:
$X=(1,2,3,6)$ has $\kappa=5$ and
$\Delta E(\kappa+3)=1<2$.
This is an auxiliary source-profile obstruction, not an original
unimodality counterexample.

\section{A general CI backward-ratio payment and the sharp gap-six refinement}
Retain the preceding source-birth theorem's notation:
$D=a-\alpha-h$, $r=a-j$, $e_0=D-2r-2$, $C=d_{r+1}-d_r>0$,
$\delta_0=\kappa_Y-\alpha$, and $s=\delta_0-e_0-2$.
Suppose $s>0$, $r\ge s-1$.
The skeleton-only primitive relative envelope and CI degree ratio
give, for $0\le i\le s-2$,
\[
              (r-i)d_{r-i}\le(r+e_0+3+i)d_{r-i-1}.     \tag{2}
\]
Consequently the FIRST birth alone pays $C$ whenever
\[
 (e_0+1)\prod_{i=0}^{s-2}(r+e_0+3+i)
       \le(r+1)\prod_{i=0}^{s-2}(r-i).                \tag{3}
\]
Indeed (2), together with $(r+1)C\le(e_0+1)d_r$,
gives $C\le d_{r-s+1}\le q_Z(r-s+1)$.
The latter reflected coefficient is at most the compulsory
birth payment, exactly as in the preceding theorem.
Empty products handle $s=1$. This is an arbitrary-gap sufficient
criterion, not a claim that (3) always holds.

\begin{theorem}[Sharp gap six for owners with at least four non-one parts]
If $h\ge4$, $F_Z(j+1)\le F_Z(j)$, the opposite SOURCE prefix
$E_Y$ is nondecreasing through $j+1$, and
$\alpha\ge\kappa_Y-6$, then $e_{n+1}\le e_n$.
With the usual OLD A2 head input, S holds at EVERY head.
The conclusion has no C or first-half premise, and the constant
six is optimal uniformly on this owner-count range.
\end{theorem}
\begin{proof}
The preceding source-gap proof handles $s\le0$, negative $r$,
and nonpositive curvature unchanged; the $r=-1$ birth is still
present since $D+h+2-\delta_0\ge2h-4\ge4$.
For positive recovery, $s\in\{1,2,3,4\}$,
$0\le e_0\le3$, $e_0+s+2=\delta_0\le6$.
For $r\ge7$, (3) follows from the following COMPLETE table:
\[
\begin{array}{c|c|l}
s& e_0\text{ at most}&\text{sufficient numerator/denominator comparison}\\\hline
1&3&4\le r+1\\
2&2&3(r+5)\le(r+1)r\\
3&1&2(r+4)(r+5)\le(r+1)r(r-1)\\
4&0&(r+3)(r+4)(r+5)\le(r+1)r(r-1)(r-2).
\end{array}
\]
The conjunction is kernel checked for ALL real $r\ge7$.
Only $0\le r\le6$ remains. Enumerate all
$D=2r+2+e_0$, $0\le e_0\le3$,
$e_0+3\le\delta_0\le6$, all partitions of $D$ into $h\ge4$
positive capacity increments, and ALL initial runs represented
by the same exact clamp $1\le\alpha\le\max(D-r+h+1-\delta_0,r+1)$.
All 33447 representatives, including 25325 positive-curvature
cases, satisfy the birth payment. Independent full-forward
reconstruction checks every case and all seven layer hashes.
This is the entire finite complement of the unbounded table.
\end{proof}

Sharpness is actual and unbounded. For $\alpha\ge2$, put
$Z=1^\alpha2^4$, $a=\alpha+8$, $j=a-1$,
$Y=1(\alpha+6)^{\alpha+7}$. Then $\kappa_Y=\alpha+7$,
$r=1$, $d_1=4$, $d_2=6$, and the local source gives
$e_n=6$, $e_{n+1}=7$. Thus gap seven cannot replace six.
The explicit S head one still pays this recovery; no original
or correctly conditioned S counterexample is asserted.

\section{Every opposite first missing weight at most eight is paid}
\begin{theorem}[Absolute first-source threshold eight, arbitrary capacities]
Under the same mixed/source-prefix assumptions and $h\ge4$,
EVERY $\kappa_Y\le8$ gives tail nonrecovery, with unrestricted
owner initial ones, opposite count and maximum capacities.
This threshold eight is optimal for the pointwise statement.
\end{theorem}
\begin{proof}
Use the exact initial exchange to choose the equivalent owner
with $\alpha_* =\max(\alpha,t_1-1)$.
The gap-six theorem handles $\kappa_Y\le7$, and also handles
$\kappa_Y=8$ whenever $\alpha_*\ge2$.
The only remaining canonical case is
\[
                    \kappa_Y=8,\quad\alpha=1,\quad t_1=2.
\]
Here $\delta_0=7$, $s\le5$, $e_0\le4$.
For $r\ge10$, the five sufficient instances of (3) are
\[
\begin{split}
5&\le r+1,\\
4(r+6)&\le(r+1)r,\\
3(r+5)(r+6)&\le(r+1)r(r-1),\\
2(r+4)(r+5)(r+6)&\le(r+1)r(r-1)(r-2),\\
(r+3)(r+4)(r+5)(r+6)&\le(r+1)r(r-1)(r-2)(r-3).
\end{split}
\]
Their complete unbounded conjunction is kernel checked.
For $r\le9$, (1) provides the TWO-term source bound
\[
              \Delta Q_{j+1}\ge A_{b_0}+A_{b_0-2},
                         \quad b_0=D-r+h-6.           \tag{4}
\]
Enumerate $D=2r+2+e_0$, $e_0=0,\ldots,4$, all $h\ge4$
partitions with $t_1=2$, and fixed $\alpha=1$.
Every one of the 10101 representatives satisfies (4)'s payment
when its primitive curvature is positive: 2758 critical cases,
three equalities. Independent full-forward enumeration checks
all ten layer hashes. These ranges exhaust the residual domain.
Negative reverse indices use the same positive birth argument,
since $h\ge4$ still makes $D+h+2-7\ge3$.
Thus all cases are paid.
\end{proof}

For sharpness take $Z=(1,2,2,2,2)$, $Y=1\,8^9$, $j=9$.
The opposite has $\kappa_Y=9$ and a nondecreasing source prefix,
while the reverse product tail recovers from three to five.
Its correctly conditioned head one pays the recovery. The binary
owner's original core family is independently proved unimodal.

These are source-FIRST-MISSING bounds, not maximum-capacity
bounds. No original least-counterexample first-weight bound is
increased merely from this scoped S result. General both-rising
late overlap and mixed cases outside the source-prefix/gap/mass
regions remain OPEN, as do general C/A2 and original undilated UI.

Six additional scoped kernel reports certify both complete large-r
polynomial tables and the exact source-step curvature payments.
The normalized-source grammar and existing early dyadic profiles
give (1); the actual scalar-to-composition assembly and two full
finite programs are written and independently checked. No end-to-end
Lean verification of those finite programs is claimed.
